\documentclass[12pt,a4paper]{ctexart}

\usepackage[a4paper,margin=2.5cm]{geometry}
\usepackage{amsmath,amssymb,amsthm,mathtools,bm}
\usepackage{mathrsfs}
\usepackage{hyperref}
\usepackage{enumitem}
\usepackage{booktabs}
\usepackage{array}
\usepackage{longtable}
\usepackage{tikz-cd}
\usepackage{fancyhdr}

\hypersetup{
    colorlinks=true,
    linkcolor=blue,
    citecolor=blue,
    urlcolor=blue
}

\pagestyle{fancy}
\fancyhf{}
\fancyhead[L]{LWE 与最坏情况格问题归约复习}
\fancyhead[R]{\thepage}

\setlength{\parskip}{0.4em}
\setlength{\parindent}{2em}

\newtheorem{definition}{定义}[section]
\newtheorem{proposition}[definition]{命题}
\newtheorem{lemma}[definition]{引理}
\newtheorem{theorem}[definition]{定理}
\newtheorem{remark}[definition]{注}

\newcommand{\R}{\mathbb{R}}
\newcommand{\Z}{\mathbb{Z}}
\newcommand{\T}{\mathbb{T}}
\newcommand{\E}{\mathbb{E}}
\newcommand{\eps}{\varepsilon}
\newcommand{\vol}{\mathrm{vol}}
\newcommand{\dist}{\mathrm{dist}}
\newcommand{\supp}{\mathrm{supp}}
\newcommand{\dual}{^{*}}
\newcommand{\ip}[2]{\left\langle #1,#2\right\rangle}
\newcommand{\norm}[1]{\left\lVert #1\right\rVert}
\newcommand{\qedherecustom}{\hfill$\square$}

\title{LWE 与最坏情况格问题归约复习提纲\\
\large 以 \texorpdfstring{$f_s(x)\approx \mathrm{Uniform}+\mathrm{Fourier\ bias}$}{fs(x) approx Uniform + Fourier bias} 与量化推导为主线}
\author{}
\date{}

\begin{document}

\maketitle
\tableofcontents
\newpage

\section{归约命题与整体框架}

\subsection{归约目标（Regev 2005）}

LWE（Learning With Errors）困难性的核心理论基础是如下归约命题：

\begin{quote}
若存在一个多项式时间算法能够解决平均情况的 LWE（判定版或搜索版），则存在一个量子多项式时间算法能够解决某些最坏情况格问题（例如 GapSVP、SIVP 的近似版本）。
\end{quote}

形式化写作
\[
\text{Worst-case lattice problem}
\;\leq_{\text{quantum}}\;
\text{Average-case LWE}.
\]

这说明 LWE 的安全性并非仅来自“随机实例难”，而是与最坏情况格问题的困难性相联系。

\subsection{归约主线结构}

本文以如下主线组织复习内容：
\[
\text{格 }L
\;\longrightarrow\;
\text{连续模格高斯 }D_{\R^n/L,s}
\;\longrightarrow\;
\text{Poisson/Fourier 展开}
\]
\[
\longrightarrow\;
f_s(x)=\text{Uniform}+\text{Fourier bias}
\;\longrightarrow\;
a=\lfloor qx\rfloor \bmod q
\;\longrightarrow\;
\text{离散 Fourier 结构}
\;\longrightarrow\;
\text{LWE}.
\]

其中最关键的两步是：

\begin{enumerate}[label=\arabic*.]
    \item 从 Poisson 求和推出
    \[
    f_s(x)\approx \text{Uniform}+\text{Fourier bias};
    \]
    \item 证明量化（rounding）之后，连续高斯模格分布不再以“高斯形状”出现，而是变成
    \[
    \text{Uniform}+\text{discrete Fourier bias}.
    \]
\end{enumerate}

\section{基本对象与记号}

\subsection{格与对偶格}

\begin{definition}[格]
设 \(b_1,\dots,b_n\in \R^n\) 线性无关，则由它们生成的格定义为
\[
L=\left\{\sum_{i=1}^n z_i b_i: z_i\in \Z\right\}.
\]
\end{definition}

\begin{definition}[对偶格]
格 \(L\subset \R^n\) 的对偶格定义为
\[
L\dual=\{w\in \R^n:\ip{w}{y}\in \Z,\ \forall y\in L\}.
\]
\end{definition}

记格的基本域体积为 \(\det L\)。

\subsection{高斯函数与分布}

定义连续高斯核
\[
\rho_s(x)=e^{-\pi \norm{x}^2/s^2},\qquad x\in \R^n.
\]
当 \(A\subseteq \R^n\) 为离散集合时，记
\[
\rho_s(A)=\sum_{x\in A}\rho_s(x).
\]

离散高斯分布 \(D_{L,s}\) 定义为
\[
\Pr[Y=y]=\frac{\rho_s(y)}{\rho_s(L)},\qquad y\in L.
\]

\subsection{LWE 分布}

给定模数 \(q\)、维度 \(n\)、秘密 \(s\in \Z_q^n\) 和误差分布 \(\chi\)，LWE 样本定义为
\[
(a,b)\in \Z_q^n\times \Z_q,\qquad
b=\ip{a}{s}+e \pmod q,\quad e\sim \chi,
\]
其中 \(a\) 通常均匀分布于 \(\Z_q^n\)。

\section{连续模格高斯与 Poisson 展开}

\subsection{连续模格高斯密度定义}

定义在商空间 \(\R^n/L\)（也可视为某个基本域上）的连续模格高斯密度为
\[
f_s(x)=\frac{1}{s^n}\sum_{y\in L}\rho_s(x+y),\qquad x\in \R^n/L.
\]
这里 \(x\) 的实际代表元并不重要，重要的是等价类 \(x\bmod L\)。

\subsection{采样视角与等价表示}

若先采样
\[
z\sim D_{\R^n,s}
\]
再取
\[
x=z\bmod L,
\]
则所得 \(x\) 的分布正是 \(D_{\R^n/L,s}\)，其密度即为 \(f_s(x)\)。

\subsection{Poisson 求和公式}

我们将使用如下形式的 Poisson 求和公式：
\[
\sum_{y\in L} g(x+y)=\frac{1}{\det L}\sum_{w\in L\dual}\widehat{g}(w)e^{2\pi i\ip{w}{x}},
\]
其中 \(\widehat{g}\) 是 \(g\) 的 Fourier 变换。

对 \(g(x)=\rho_s(x)\) 而言，有
\[
\widehat{\rho_s}(w)=s^n \rho_{1/s}(w).
\]
因此得到
\[
\sum_{y\in L}\rho_s(x+y)
=
\frac{s^n}{\det L}\sum_{w\in L\dual}\rho_{1/s}(w)e^{2\pi i\ip{w}{x}}.
\]
两边同除以 \(s^n\)，便得
\[
f_s(x)
=
\frac{1}{\det L}\sum_{w\in L\dual}\rho_{1/s}(w)e^{2\pi i\ip{w}{x}}.
\]
这就是后续所有分析的起点。

\subsection{Poisson 公式在本问题中的单独推导}

本节单独推导上式，因为它是整个复习中的关键桥梁。

\subsubsection*{第一步：从一般 Poisson 求和出发}

对足够良好的函数 \(g:\R^n\to \mathbb{C}\)，Poisson 求和公式给出
\[
\sum_{y\in L} g(x+y)
=
\frac{1}{\det L}\sum_{w\in L\dual}\widehat{g}(w)e^{2\pi i\ip{w}{x}}.
\]
这里 Fourier 变换按如下约定：
\[
\widehat{g}(w)=\int_{\R^n} g(t)e^{-2\pi i\ip{w}{t}}\,dt.
\]

\subsubsection*{第二步：取 \(g=\rho_s\)}

取
\[
g(t)=\rho_s(t)=e^{-\pi \norm{t}^2/s^2}.
\]
标准高斯 Fourier 变换公式告诉我们
\[
\widehat{\rho_s}(w)=s^n e^{-\pi s^2 \norm{w}^2}=s^n \rho_{1/s}(w).
\]

\subsubsection*{第三步：代回 Poisson 公式}

于是
\[
\sum_{y\in L}\rho_s(x+y)
=
\frac{1}{\det L}
\sum_{w\in L\dual}
\left(s^n\rho_{1/s}(w)\right)e^{2\pi i\ip{w}{x}}.
\]
即
\[
\sum_{y\in L}\rho_s(x+y)
=
\frac{s^n}{\det L}\sum_{w\in L\dual}\rho_{1/s}(w)e^{2\pi i\ip{w}{x}}.
\]

\subsubsection*{第四步：恢复 \(f_s(x)\)}

由于
\[
f_s(x)=\frac{1}{s^n}\sum_{y\in L}\rho_s(x+y),
\]
故
\[
f_s(x)
=
\frac{1}{\det L}\sum_{w\in L\dual}\rho_{1/s}(w)e^{2\pi i\ip{w}{x}}.
\]
推导完成。

\paragraph{说明}
这个公式的意义在于：连续模格高斯分布在位置空间中的周期结构，可以完全转写为对偶格上的 Fourier 模式之和。也正因此，对偶格中的短向量会在分布中留下“低频偏移”。

\section{核心结论一：\texorpdfstring{$f_s(x)=\mathrm{Uniform}+\mathrm{Fourier\ bias}$}{fs(x)=Uniform+Fourier bias}}

本节是全文第一条主线。

\subsection{分离零频率项}

由
\[
f_s(x)
=
\frac{1}{\det L}\sum_{w\in L\dual}\rho_{1/s}(w)e^{2\pi i\ip{w}{x}},
\]
将 \(w=0\) 项单独提出：
\[
f_s(x)
=
\frac{1}{\det L}\rho_{1/s}(0)
+
\frac{1}{\det L}\sum_{w\in L\dual\setminus\{0\}}
\rho_{1/s}(w)e^{2\pi i\ip{w}{x}}.
\]
因为 \(\rho_{1/s}(0)=1\)，故
\[
f_s(x)
=
\frac{1}{\det L}
+
\frac{1}{\det L}\sum_{w\in L\dual\setminus\{0\}}
\rho_{1/s}(w)e^{2\pi i\ip{w}{x}}.
\]

\subsection{均匀项的解释}

常数项
\[
\frac{1}{\det L}
\]
恰是基本域上的均匀密度，因为基本域体积即为 \(\det L\)。

因此：
\[
\text{主项}=\text{Uniform density on }\R^n/L.
\]

\subsection{Fourier 偏移项的解释}

剩余项
\[
\Delta(x)=
\frac{1}{\det L}\sum_{w\in L\dual\setminus\{0\}}
\rho_{1/s}(w)e^{2\pi i\ip{w}{x}}
\]
即为非零频率模式的叠加。它表示分布相较于完全均匀分布的偏离，因此称之为 Fourier bias。

\subsection{结论形式}

所以严格地说，
\[
f_s(x)
=
\underbrace{\frac{1}{\det L}}_{\text{Uniform}}
+
\underbrace{
\frac{1}{\det L}\sum_{w\in L\dual\setminus\{0\}}
\rho_{1/s}(w)e^{2\pi i\ip{w}{x}}
}_{\text{Fourier bias}}.
\]
这就是
\[
f_s(x)=\mathrm{Uniform}+\mathrm{Fourier\ bias}
\]
的精确来源：它不是启发式说法，而是 Poisson 展开中将零频率项与非零频率项分离后的严格表达。

\section{平滑参数（Smoothing Parameter）复习}

本节单列复习，因为它是控制 Fourier bias 总量、保证近似均匀性的核心工具。

\subsection{定义}

\begin{definition}[平滑参数]
给定格 \(L\subset \R^n\) 和 \(\eps>0\)，其平滑参数 \(\eta_\eps(L)\) 定义为满足
\[
\rho_{1/s}(L\dual\setminus\{0\})=\sum_{w\in L\dual\setminus\{0\}} \rho_{1/s}(w)\le \eps
\]
的最小 \(s\)（精确地说是下确界）：
\[
\eta_\eps(L)
=
\inf\left\{
s>0:
\sum_{w\in L\dual\setminus\{0\}}\rho_{1/s}(w)\le \eps
\right\}.
\]
\end{definition}

\subsection{基本解释}

平滑参数衡量的是：当高斯参数 \(s\) 足够大时，对偶格上的非零频率总质量会被压得多小。

换言之，当 \(s\ge \eta_\eps(L)\) 时，对偶格上的非零 Fourier 模式总贡献不超过 \(\eps\)。

\subsection{基本性质}

\begin{enumerate}[label=\arabic*.]
    \item \textbf{单调性：} 若 \(s_1\le s_2\)，则
    \[
    \rho_{1/s_2}(L\dual\setminus\{0\})\le \rho_{1/s_1}(L\dual\setminus\{0\}).
    \]
    \item \textbf{尺度变换：} 对任意 \(\alpha>0\)，有
    \[
    \eta_\eps(\alpha L)=\alpha \eta_\eps(L).
    \]
    \item \textbf{与对偶格短向量的关系：} 若 \(L\dual\) 存在很短的非零向量，则需要更大的 \(s\) 才能把 Fourier 尾部压小。
\end{enumerate}

\subsection{与近似均匀性的关系}

若 \(s\ge \eta_\eps(L)\)，则对任意 \(x\in \R^n/L\)，有
\[
\left|
f_s(x)-\frac{1}{\det L}
\right|
\le
\frac{1}{\det L}\sum_{w\in L\dual\setminus\{0\}}\rho_{1/s}(w)
\le
\frac{\eps}{\det L}.
\]
因此
\[
f_s(x)\approx \frac{1}{\det L},
\]
即连续模格高斯分布接近基本域上的均匀分布。

\subsection{本归约中的作用}

平滑参数的作用可以总结为：

\begin{enumerate}[label=\arabic*.]
    \item 控制非零 Fourier 项总量；
    \item 保证 \(f_s(x)\) 接近均匀；
    \item 为后续量化后的 \(a\) 近似均匀提供基础。
\end{enumerate}

\section{核心结论二：量化后的离散分布推导}

本节是全文第二条主线。目标是解释：为什么经过 rounding/quantization 之后，原本的连续模格高斯分布不再以“高斯形状”出现，而是变成离散 Fourier 结构。

\subsection{量化映射定义}

设
\[
a=\lfloor qx\rfloor \bmod q.
\]
在多维情形中，这是逐坐标定义的，故 \(a\in \Z_q^n\)。

\subsection{一维情况的严格推导}

先考虑一维，设 \(X\) 是定义在 \([0,1)\) 上的随机变量，密度为 \(f(x)\)，定义
\[
A=\lfloor qX\rfloor,\qquad A\in \{0,1,\dots,q-1\}.
\]
则
\[
\Pr[A=r]
=
\int_{r/q}^{(r+1)/q} f(x)\,dx.
\]
这是因为事件 \(A=r\) 恰对应 \(X\in [r/q,(r+1)/q)\)。

\subsection{代入 Fourier 展开}

若
\[
f(x)=1+\sum_{t\neq 0}\widehat{f}(t)e^{2\pi i t x},
\]
则
\[
\Pr[A=r]
=
\int_{r/q}^{(r+1)/q}
\left(
1+\sum_{t\neq 0}\widehat{f}(t)e^{2\pi i t x}
\right)\,dx.
\]
拆分后得
\[
\Pr[A=r]
=
\frac{1}{q}
+
\sum_{t\neq 0}\widehat{f}(t)
\int_{r/q}^{(r+1)/q}e^{2\pi i t x}\,dx.
\]

\subsection{计算指数项积分}

对任意 \(t\neq 0\)：
\[
\int_{r/q}^{(r+1)/q}e^{2\pi i t x}\,dx
=
e^{2\pi i tr/q}\int_0^{1/q}e^{2\pi i t u}\,du
=
e^{2\pi i tr/q}\cdot \frac{e^{2\pi i t/q}-1}{2\pi i t}.
\]
故
\[
\Pr[A=r]
=
\frac{1}{q}
+
\sum_{t\neq 0}
\widehat{f}(t)\,
e^{2\pi i tr/q}\,
\frac{e^{2\pi i t/q}-1}{2\pi i t}.
\]

\subsection{结论：离散 Fourier 结构}

上式说明，量化后的分布不是“仍然保持高斯形状”，而是
\[
\Pr[A=r]
=
\underbrace{\frac{1}{q}}_{\text{Uniform}}
+
\underbrace{
\sum_{t\neq 0}
\widehat{f}(t)\,
e^{2\pi i tr/q}\,
\frac{e^{2\pi i t/q}-1}{2\pi i t}
}_{\text{discrete Fourier bias}}.
\]

\subsection{多维情况}

若
\[
a=\lfloor qx\rfloor \bmod q,\qquad a\in \Z_q^n,
\]
则对每个 \(a\in \Z_q^n\) 有
\[
\Pr[A=a]
=
\int_{a/q+[0,1/q)^n} f_s(x)\,dx.
\]
若
\[
f_s(x)=\frac{1}{\det L}+\sum_{w\neq 0} c_w e^{2\pi i\ip{w}{x}},
\]
代入可得
\[
\Pr[A=a]
=
\frac{1}{q^n}
+
\sum_{w\neq 0}
c_w\,
e^{2\pi i\ip{w}{a}/q}\,
K_q(w),
\]
其中
\[
K_q(w)=
\int_{[0,1/q)^n}
e^{2\pi i\ip{w}{u}}\,du.
\]
因此多维下同样得到
\[
\text{quantized distribution}
=
\text{Uniform}
+
\text{discrete Fourier bias}.
\]

\subsection{为什么说 rounding “打破”了高斯结构}

准确地说，rounding 并不是“把信息消掉”，而是将连续密度通过一个非线性的、分箱式（binning）的推前映射
\[
x\mapsto \lfloor qx\rfloor \bmod q
\]
压缩成有限离散分布。高斯分布对线性变换封闭，但对这种非线性的 floor/round 操作并不封闭，因此量化后保留下来的结构不再是“钟形曲线”，而是每个小箱子上的概率质量。对于 periodic Gaussian，这些概率质量自然以离散 Fourier 形式表达。

\section{从 Fourier bias 到 LWE 分布}

在归约中，量化后得到的离散 Fourier 偏差会进一步转写为 LWE 形式。直观上：

\begin{itemize}
    \item 连续空间中的频率项
    \[
    e^{2\pi i\ip{w}{x}}
    \]
    经过量化后，变成关于 \(a\in \Z_q^n\) 的离散相位
    \[
    e^{2\pi i\ip{w}{a}/q}.
    \]
    \item 当这种相位结构与高斯平滑结合时，可以等价地写成带噪声的线性关系
    \[
    b=\ip{a}{s}+e \pmod q.
    \]
\end{itemize}

因此，LWE 的本质可以概括为：

\begin{quote}
线性相位结构在模 \(q\) 离散化和高斯模糊之后，表现为“线性函数 + 小噪声”的样本。
\end{quote}

\section{利用 LWE 判定器完成归约}

若存在一个算法可以区分：

\begin{enumerate}[label=\arabic*.]
    \item 完全均匀分布；
    \item 含有某个隐藏线性频率结构的分布；
\end{enumerate}
则该算法就可以检测是否存在显著的对偶格频率项，也即是否存在某些短对偶向量。进一步结合归约构造，可以把这种检测能力转化为对某些最坏情况格问题（如 GapSVP、SIVP 近似版）的求解能力。

\section{关键概念总结}

\begin{enumerate}[label=\arabic*.]
    \item \textbf{LWE 的本质：} 线性相位 + 高斯模糊。
    \item \textbf{\(a\) 的分布：} 近似均匀，而不是高斯。
    \item \textbf{bias 来源：} 对偶格中的短向量。
    \item \textbf{\(b\) 的本质：} 对隐藏 Fourier 结构的观测变量。
    \item \textbf{噪声来源：} Fourier Gaussian 与离散化共同作用。
    \item \textbf{\(x\) 的角色：} 连接“连续模格高斯”与“离散 LWE 样本”的桥梁变量；只需关心其模格等价类。
\end{enumerate}

\section{复习重点建议}

本文复习的优先级建议如下：

\begin{enumerate}[label=\arabic*.]
    \item 首先掌握 Poisson 公式在本问题中的具体推导；
    \item 其次掌握
    \[
    f_s(x)=\mathrm{Uniform}+\mathrm{Fourier\ bias}
    \]
    的严格分解；
    \item 再掌握量化后离散分布的推前公式与 discrete Fourier bias；
    \item 单独复习平滑参数及其性质；
    \item 最后再回到 LWE 形式本身，理解为什么这些 Fourier 结构会转写为
    \[
    b=\ip{a}{s}+e \pmod q.
    \]
\end{enumerate}

\end{document}